\documentclass[11pt]{article}
\usepackage{hyphenat}
\usepackage{natbib}
\usepackage[T1]{fontenc}
\usepackage[utf8]{inputenc}
\usepackage{lmodern}
\usepackage{microtype}
\usepackage{amsmath,amssymb,amsfonts}
\usepackage{booktabs}
\usepackage{array}
\usepackage{graphicx}
\usepackage{xcolor}
\usepackage{setspace}
\usepackage{tcolorbox}
\tcbuselibrary{skins,breakable}

\usepackage{hyperref}
\hypersetup{
    colorlinks=true,
    linkcolor=blue!80!black,
    citecolor=blue!80!black,
    urlcolor=blue!80!black
}

\definecolor{revgray}{rgb}{0.96, 0.96, 0.98}
\definecolor{revborder}{rgb}{0.75, 0.78, 0.85}

\newtcolorbox{reviewercomment}[1][]{
    enhanced,
    breakable,
    colback=revgray,
    colframe=revborder,
    coltitle=black,
    arc=2mm,
    boxrule=0.8pt,
    left=12pt,
    right=12pt,
    top=9pt,
    bottom=9pt,
    fontupper=\itshape,
    before skip=10pt,
    after skip=10pt,
    #1
}

\setlength{\parindent}{0pt}
\setlength{\parskip}{7pt}
\setstretch{1.12}

\begin{document}

\noindent{\LARGE \textbf{Response to Reviewers and Editors}}\\[6pt]
\textbf{Manuscript Number:} SPQ-26-0082\\
\textbf{Date:} \today\\[12pt]

\section*{Overview and Letter to the Editors}

Dear Professors Slattery Walker and Dippong,

Thank you for the opportunity to revise and resubmit my manuscript, \textit{Patriotism as a Choice Process: National Pride After 9/11} (SPQ-26-0082), for publication in \textit{Social Psychology Quarterly}. I am grateful to you and the two reviewers for the careful, constructive reading of the paper. Their comments have pushed me to spell out the theoretical mechanisms underlying my two competing perspectives in much greater detail, to justify and expand my empirical and modeling strategy, and to substantially polish the manuscript's presentation and flow.

In response to this feedback, I have undertaken a comprehensive revision of the manuscript. Below, I provide a point-by-point response to every issue raised by the editors and the two reviewers. In summary, my primary revisions include: (1) substantially expanding the theoretical mechanism sections to explain why 9/11 should weaken the proximal attribution rule and strengthen the distal rule, and why saliency should be stronger for individuals embedded in nested subgroups, grounding both claims in Lawler's (1992) own boundary conditions and in \citet{lawler2006commitment-b3a}; (2) adding a dedicated discussion of the 9/11 event itself, its national-scale character, and its relation to prior empirical work on 9/11 and comparable threat events \citep{perrin2005national, schatz2007waving, sullivan1992patriotism}; (3) adding a full Analytic Strategy subsection justifying my modeling choices, including a heteroskedasticity-robust standard-error check; (4) completing the introduction and moving the descriptive statistics table into the main text; (5) integrating the treatment of Social Identity Theory into the main text and clarifying its relationship to the authoritarian-personality tradition, per Reviewer~2; (6) introducing the proximal/distal distinction earlier, in the Introduction; (7) a general copyediting and polish pass throughout; (8) removing the supplemental political-ideology appendix in favor of a footnote, per Reviewer~2; and (9) replacing the undefined term ``excessive'' attachment with Schatz, Staub, and Lavine's (1999) qualitative ``blind patriotism'' construct, per Reviewer~1.

Below, I detail each editor and reviewer comment verbatim, followed by the precise location in the revised manuscript and my substantive response.

\clearpage

\section*{Response to the Editors}

\subsection*{Point 1: Theoretical Mechanisms Are Underdeveloped}

\begin{reviewercomment}
Both reviewers see much of value in the paper. However, both feel that the work is undertheorized; that is, the mechanisms underlying the theoretical processes are not sufficiently spelled out.
\end{reviewercomment}

\noindent\textbf{Location in Revised Manuscript:} Section~2.2 (Lawler's Choice-Process Theory and Nested Attachments) and Section~2.3 (The Theoretical Status of 9/11: A Shift in Attribution Rules).\\[4pt]
\textbf{Response:} I have substantially expanded the mechanism underlying the choice-process account. The revised Section~2.2 now specifies, following Lawler's (1992, pp.~610--620) own boundary conditions, exactly when the default proximal rule reverses into a distal rule: when the rituals and symbols of the larger collectivity become salient enough to override the interactional advantage ordinarily enjoyed by proximal subgroups, or when structural circumstances make the larger unit the actual source of an individual's sense of control. I further draw on \citet{lawler2006commitment-b3a}, whose experimental distinction between structurally enabled and structurally induced exchange relations supplies the individual-level mechanism (a stronger sense of control, more positive emotion, and more durable commitment when actors can act on their preferred affiliations) that I extend across levels of nesting. Section~2.3 then applies this explicitly to 9/11: rather than asserting an unexplained ``shock,'' I argue that the extraordinary volume of national rituals and symbols in its aftermath (flag displays, moments of silence, memorial services, saturation media coverage) is precisely the condition Lawler's theory identifies as capable of triggering the shift from a proximal to a distal attribution rule. I also address the reviewer's specific query about why saliency should be stronger for individuals embedded in nested subgroups (see my response to Reviewer 1, Point 1 below for the parallel revision).

\subsection*{Point 2: Presentation and Flow}

\begin{reviewercomment}
In addition, both reviewers think that the paper needs to be a little more polished in its presentation and that the flow of the paper could be improved.
\end{reviewercomment}

\noindent\textbf{Location in Revised Manuscript:} Throughout.\\[4pt]
\textbf{Response:} This point is the same underlying concern raised by Reviewer 2 in their Point 3 below; please see my response there. In brief, I revised the prose throughout the manuscript for flow and clarity, tightening the theory sections, the introduction of the proximal/distal distinction, the Data and Variables section, the Analytic Strategy subsection, and the Discussion as I addressed the editors' and reviewers' other comments, and I subsequently went through the remainder of the manuscript to shorten overly long sentences and clarify paragraph transitions.

\subsection*{Point 3: Methods Section Detail and Justification}

\begin{reviewercomment}
Additionally, one reviewer believes, and we agree, that more detail and justification is needed in the methods section.
\end{reviewercomment}

\noindent\textbf{Location in Revised Manuscript:} New Section~3.5 (Analytic Strategy).\\[4pt]
\textbf{Response:} This is the same underlying concern raised by Reviewer~1 in their Point~3(c) below; please see my fuller response there. In brief, I added a new \textit{Analytic Strategy} subsection that specifies my estimator (weighted OLS, using the GSS full-sample probability weight \texttt{wtssall}), distinguishes the two model designs used in the paper (pooled 1996--2004 models with a year~$\times$~education interaction, versus 2004-only cross-sectional models for the voluntary-association analysis), and describes my handling of missing data at the scale-construction level. I also added a heteroskedasticity-robust (HC1) standard-error robustness check for the three focal coefficients underlying my hypothesis tests, reported narratively in this subsection, to address the broader concern about methodological transparency.

\clearpage

\section*{Response to Reviewer 1}

\subsection*{Point 1: Theoretical Mechanisms, Outcome Variables, and Unsupported Expectations}

\begin{reviewercomment}
First, the theoretical arguments lack clear substance. While the overarching theories are briefly referred to, the underlying mechanisms by which 9/11 should spark more/less nationalism and pride among different subgroups remain unclear. Moreover, too little attention is paid to social-psychological differences between the two outcome variables. What is the benefit of studying both of them? How do they differ? Furthermore, in several cases, the expectations are not backed up by appropriate arguments. For example, in the first paragraph on p. 5, it is unclear, why exactly we should expect ``a weakening of the proximal rule and the emergence of a distal rule.'' In the following paragraph, it is not specified, why the sense of saliency ``should be stronger for individuals embedded in nested subgroups.''
\end{reviewercomment}

\noindent\textbf{Location in Revised Manuscript:} Section~2.2 (Lawler's Choice-Process Theory and Nested Attachments), Section~2.3 (The Theoretical Status of the Post-9/11 Period: A Shift in Attribution Rules), and the new Section~3.2.1 (Why Two Outcomes? Affective Attachment versus Blind Nationalism).\\[4pt]
\textbf{Response:} I agree that the original submission asserted, rather than argued for, the shift from a proximal to a distal attribution rule and the heightened saliency of nested subgroups, without adequate theoretical support, and that it did not explain why I examine two outcome variables. I address each issue in turn.

\textit{(a) Why a shift from a proximal to a distal rule, and why should saliency be stronger for nested subgroups?} I now ground the claim that 9/11 should weaken the proximal rule directly in Lawler's (1992, pp.~610--620) own specification of when a distal rule emerges: when the rituals and symbols of the larger collectivity become salient enough to override the interactional advantage ordinarily enjoyed by proximal subgroups. I argue that the extraordinary volume of national rituals and symbols following 9/11---flag displays, moments of silence, memorial services, saturation media coverage---plausibly met this condition, rather than treating the shift as an unexplained shock. I further explain why saliency should be stronger for individuals embedded in nested subgroups by drawing on \citet{lawler2006commitment-b3a}, whose experimental distinction between structurally enabled and structurally induced exchange relations shows that the joint, shared-responsibility conditions required for control-based positive emotion to be credited to a social unit are met specifically within nested subgroups that individuals can actually act on (e.g., voluntary associations), rather than merely occupy. I explain that this same structural position supplies the raw material of positive affect that becomes available for reattribution to the nation once the distal rule is active.

\textit{(b) Why study both the U.S. Pride and Nationalism scales, and how do they differ?} This was a substantive gap in the original submission rather than a presentational one: nothing in the theory section explained why two outcomes were needed or what distinguished them conceptually. I have added a new subsubsection (Section~3.2.1) directly after the two scales are defined to make this explicit. I argue that the two scales operationalize conceptually distinct constructs rather than measuring the same thing. The U.S. Pride scale captures \textit{affective attachment} in Lawler's specific sense---a positive, emotionally invested identification with the country across substantive achievement domains, generated by choice and control experienced in everyday structured contexts and reattributed to a social unit. The Nationalism scale instead captures a comparative, ideological orientation---uncritical, ``blind'' loyalty and belief in American superiority \citep{schatz1999varieties}---closer to the construct that needs-based and authoritarian-personality accounts are theories of.

I explain that this distinction generates divergent, testable predictions rather than being a purely definitional exercise: a needs-based account implies both outcomes should move together, concentrated among the same low-status, low-autonomy groups, since it offers no principled reason for affective pride and ideological nationalism to respond differently to a common threat, whereas the choice-process account implies a more selective pattern---increased affective attachment among structurally autonomous individuals, with no corresponding shift in the more purely ideological Nationalism scale, since the choice-process mechanism is a theory of attribution for self-generated positive emotion and has no equivalent account of why beliefs about national superiority should track structural autonomy. Studying both scales jointly therefore serves as a test of discriminant validity between the two theoretical accounts. I note explicitly that my results are consistent with this second, more selective pattern (education predicts the post-9/11 increase in Pride but not in Nationalism), and I flag this dissociation as a substantive finding in its own right in the Discussion (Section~5.2).

\subsection*{Point 2: Missing Discussion of the 9/11 Event}

\begin{reviewercomment}
Second, there is virtually no discussion of the 9/11 event itself. Why is it important? How should it shape feelings of national pride? What was the context? At the moment, 9/11 appears somewhat out of the blue in the manuscript. More broadly, I miss a discussion of how the study fits in with other studies on 9/11 and other major events as well as their effects on feelings of pride and nationalism.
\end{reviewercomment}

\noindent\textbf{Location in Revised Manuscript:} New opening paragraph of Section~2.3 (The Theoretical Status of the Post-9/11 Period: A Shift in Attribution Rules).\\[4pt]
\textbf{Response:} I agree that the original submission introduced 9/11 without establishing why the event itself mattered, what its national-scale character consisted of, or how my study relates to prior empirical work on 9/11 and comparable events. I have added a new opening paragraph to Section~2.3 that addresses each of these gaps before turning to the choice-process mechanism.

The new paragraph first establishes why 9/11 is a plausible trigger for a shift in national attribution rules: the attacks were a coordinated strike on iconic national sites, covered saturation-style by national media, and experienced as an assault on the nation as a collectivity rather than on any particular locality or subgroup. I then situate this claim in prior empirical literature rather than asserting it. Drawing on a hand-coded sample of letters to the editor from before and after the attacks, \citet{perrin2005national} shows that expressions of both authoritarian and antiauthoritarian political rhetoric increased sharply in the month following September 11, indicating that the attacks activated political-cultural scripts at a national scale. \citet{schatz2007waving} independently document a parallel upsurge in concern for national symbols and ritualistic-ceremonial activity after 9/11, explicitly paralleling the similar flourish of patriotic symbolism that followed the 1991 Gulf War---directly addressing the reviewer's request for a discussion of ``other major events'' with comparable effects. I further cite \citet{sullivan1992patriotism}'s finding that a ``symbolic'' patriotism construct predicted distinct political attitudes and candidate support during the 1988 U.S.\ presidential campaign, to show that patriotic sentiment tracking the political climate of the moment is a more general phenomenon that 9/11 exemplifies rather than originates.

Having established this context, the new paragraph closes by clarifying what the present study adds beyond this existing literature: prior work establishes \textit{that} 9/11 (and comparable threat events) shifted national sentiment and discourse, but says little about \textit{which} segments of the population were most affected and \textit{why}---the distributional question that is the object of this paper, and the one on which the choice-process and needs-based accounts generate diverging, testable predictions. This connects directly into the existing discussion of the proximal-to-distal attribution shift that follows.

\subsection*{Point 3: Empirical Analysis Is Too Simplistic and Not Transparent}

\begin{reviewercomment}
Third, the empirical analysis is too simplistic and not transparent enough. The authors run logistic regression models on two waves of GSS data. There are several problems: 
\begin{itemize}
\item Many things happened between 1996 and 2004. 9/11 was certainly an important event. But there were other important events, too (e.g., war under the Bush administration), which might explain changes in pride in armed forces. It is by no means obvious that all the change observed can be attributed to 9/11.
\item Why are logistic regressions used? As I understand the data and variables section, the outcome variables are continuous scales. Why do you not use linear regressions instead?
\item The methods section that explains the modelling approach is completely missing.
\item There are no descriptive tables/plots, which renders it difficult to fully make sense of the regression results. The authors included a placeholder, but the actual table is not there.
\item Modelling choices are not justified. E.g., why are the factors underlying one outcome scale dichotomized, but not those underlying the other?
\end{itemize}
\end{reviewercomment}

\noindent\textbf{Location in Revised Manuscript:} Introduction, Section~2.3 (The Theoretical Status of the Post-9/11 Period), Section~3.2 and Section~3.5 (Data and Variables / Analytic Strategy), Section~4 (Results), the Descriptive Statistics table, and Section~5.3 (Limitations and Future Research).\\[4pt]
\textbf{Response:} I address each sub-issue in turn below.

\textit{(a) Confounding events between 1996 and 2004:} I agree that 9/11 cannot be isolated from the war in Afghanistan, the Iraq War, and the broader climate of wartime mobilization that followed it, particularly for the armed-forces pride item. Rather than attempting to partial out these events as confounds, I have reframed the paper's theoretical object of interest throughout the revised manuscript (Introduction, Section~2.3, and Discussion) as the \textit{post-9/11 period} as a whole, with the September 11 attacks serving as the catalyzing event that opened this period rather than as an isolated, punctual stimulus whose unique effect I claim to identify. Substantively, I argue that the subsequent wars are not external to the mechanism I theorize but continue to supply the same kind of salience-raising national rituals and symbols (troop deployments, wartime rhetoric, media coverage) that I argue drove the shift from a proximal to a distal attribution rule. I have added a paragraph to the Limitations section (Section~5.3) explicitly acknowledging that, because the 2004 wave falls more than two years after the attacks, I cannot cleanly decompose the contribution of 9/11 itself from the subsequent period of wartime mobilization, and I note that repeated cross-sections spanning 2001--2004 would be needed to do so.

\textit{(b) Justification for logistic vs.\ linear regression:} The reviewer's concern stems from an ambiguity in my original wording rather than from my actual modeling strategy. My two substantive dependent variables---the U.S. Pride Scale (a 0--10 additive index) and the Nationalism Scale (the mean of four items on a 1--5 scale)---are indeed treated as continuous, and every model that tests my substantive hypotheses (Tables~Table 1 (Predictors of Nationalism Scale), Table 2 (Effects of Education on U.S. Pride Scale), and Table 3 (Predictors of U.S. Pride, 2004 Sample)) is estimated via OLS linear regression, as reflected in the $R^2$, adjusted $R^2$, $F$-statistic, and RMSE diagnostics reported in each table.

The only place logistic regression appears in the paper is in my preliminary, purely descriptive analysis of the ten individual binary items that comprise the U.S. Pride Scale (e.g., ``very proud'' vs.\ not, for each domain such as armed forces, history, or the arts). Because each of these ten underlying items was coded as dichotomous, logistic regression is the appropriate choice for characterizing the item-by-item period shift illustrated in Figure 1; this analysis is meant only to document \textit{that} pride increased broadly across item domains before I turn to the composite scales. I have revised the opening of the Results section (Section~4.1) to explicitly state that these item-level logistic models are a preliminary descriptive step, distinct from the linear regression models used for all subsequent hypothesis tests, to avoid any suggestion that logistic regression was used to model the continuous scales themselves.

\textit{(c) Missing methods section:} I have added a new subsection, \textit{Analytic Strategy} (Section~3.5 of the revised manuscript), that details my estimation approach. In brief: (i)~because the U.S. Pride and Nationalism scales are continuous, all substantive hypothesis tests are estimated via weighted least squares (OLS) regression, with respondents weighted by the GSS full-sample probability weight (\texttt{wtssall}); (ii)~I distinguish two model designs---pooled 1996--2004 models with a year indicator and a year~$\times$~education interaction used to test whether the education effect changed across periods (Table 1 (Predictors of Nationalism Scale) and Table 2 (Effects of Education on U.S. Pride Scale)), versus 2004-only cross-sectional models for the nested-subgroup membership analysis (Table 3 (Predictors of U.S. Pride, 2004 Sample)), since the voluntary association battery was administered only in the 2004 wave; and (iii)~I describe my handling of missing data, noting that the U.S. Pride scale requires complete responses across all ten component items (listwise deletion at the scale level) while the Nationalism scale is computed as the mean of available items, and that analytic sample sizes are reported in each table.

\textit{(d) Missing descriptive tables/plots:} I have moved the full descriptive statistics table into the Data and Variables section (Section~3.5) and confirmed that all figures and regression tables referenced in the manuscript are included with real values rather than placeholders.

\textit{(e) Justification for dichotomization of scale items:} The two scales are dichotomized (or not) for different substantive and measurement reasons, and I have added language to Section~3.2 clarifying this choice.

The U.S. Pride scale follows the standard approach in the cross-national national-pride literature \citep{smith2006national}, which treats the response ``very proud'' as qualitatively distinct from the three lower categories (``somewhat proud,'' ``not very proud,'' ``not proud at all''). Because the theoretical quantity of interest is the number of domains in which a respondent expresses maximal, unambiguous national pride, each item is dichotomized (1 = ``very proud,'' 0 = otherwise) and the ten resulting indicators are summed into a 0--10 count. Treating the middle categories as equivalent to ``very proud'' would blur a conceptually meaningful distinction between full-throated pride and lukewarm or absent pride.

The Nationalism scale, by contrast, is built from four Likert agree/disagree items adapted from Schatz, Staub, and Lavine's (1999) work on ``blind'' versus ``constructive'' patriotism that are designed to capture graded agreement with statements of national superiority and blind loyalty; here, the degree of agreement (not simply agreement vs.\ disagreement) is the theoretically meaningful quantity, consistent with standard psychometric practice for attitude scales, so the items are reverse-coded and averaged rather than dichotomized (Cronbach's $\alpha = 0.64$).

To confirm that my substantive conclusions do not hinge on this dichotomization choice, I constructed an alternative, non-dichotomized version of the U.S. Pride scale as the mean of the ten items (reverse-coded so that higher values indicate greater pride, without collapsing categories). This continuous alternative correlates at $r = 0.87$ with my dichotomized additive scale, and re-estimating the pooled education~$\times$~year interaction model on this alternative scale yields the same substantive conclusion: the interaction remains positive (b~$=0.015$) and statistically significant under classical weighted-OLS standard errors ($p = 0.037$), and remains marginally significant under heteroskedasticity-robust standard errors ($p = 0.073$, one-tailed $p = 0.037$). I now report this robustness check in a footnote attached to the description of the U.S. Pride scale in Section~3.2 of the revised manuscript.

\medskip
\noindent\textbf{Robustness check: heteroskedasticity-robust standard errors.} In response to the editors' and Reviewer~1's requests for greater methodological transparency and justification, I also re-estimated all three of my focal hypothesis-testing models using heteroskedasticity-consistent (HC1) standard errors, as a complement to the classical weighted-OLS standard errors reported in the manuscript's main tables. (Because my data are a pooled cross-section of independent GSS respondents rather than a clustered or multi-stage complex sample as fielded, and because the GSS full-probability weight \texttt{wtssall} is already used to weight all models, heteroskedasticity-robust rather than cluster-robust standard errors are the relevant design-based correction here.) Table~\ref{tab:robustse} summarizes the results for the three coefficients central to my argument.

\begin{table}[htbp]
\centering
\caption{Classical Weighted-OLS vs.\ Heteroskedasticity-Robust (HC1) Standard Errors for Focal Coefficients}
\label{tab:robustse}
\resizebox{\textwidth}{!}{%
\begin{tabular}{>{\raggedright\arraybackslash}p{4.6cm} >{\raggedright\arraybackslash}p{3.2cm} r r r c}
\toprule
\textbf{Model} & \textbf{Coefficient} & \textbf{Estimate} & \textbf{SE (classical)} & \textbf{SE (HC1)} & \textbf{$p$ (classical / HC1)} \\
\midrule
Nationalism, Educ.\ $\times$ Year (Table 1) & Educ.\ $\times$ Year 2004 & $-0.003$ & 0.010 & 0.011 & 0.789 / 0.810 \\
Pride, Educ.\ $\times$ Year (Table 2) & Educ.\ $\times$ Year 2004 & 0.114 & 0.044 & 0.050 & 0.010 / 0.021 \\
Full 2004 model (Table 3) & Total memberships & 0.112 & 0.043 & 0.051 & 0.010 / 0.027 \\
\bottomrule
\end{tabular}%
}
\end{table}

As Table~\ref{tab:robustse} shows, robust standard errors are modestly larger than the classical weighted-OLS standard errors (roughly 8--15\% larger across coefficients), as would be expected once heteroskedasticity induced by the survey weights is accounted for. Critically, none of my substantive conclusions change: the positive education~$\times$~year interaction on the U.S. Pride scale remains statistically significant ($p = 0.021$ under HC1), the corresponding interaction on the Nationalism scale remains clearly non-significant under both specifications (consistent with my claim that the choice-process story, not the irrational-patriotism story, is supported), and the effect of total voluntary association memberships in the full 2004 model remains statistically significant ($p = 0.027$ under HC1). I separately note that the simple 1996-only education coefficient (Table 2 (Effects of Education on U.S. Pride Scale), Column~1) moves from $p = 0.030$ under classical standard errors to a borderline $p = 0.055$ under HC1 (though it remains significant one-tailed, $p = 0.028$, consistent with my directional hypothesis); I report this nuance transparently in the text rather than treating the 1996 effect as unambiguous. I now report this robustness check narratively in the newly added Analytic Strategy subsection (Section~3.5) of the revised manuscript, rather than adding separate robust-SE columns to Tables 1--3, in the interest of keeping the main tables uncluttered; the full comparison in Table~\ref{tab:robustse} above is available to the editors and reviewers on request.

\subsection*{Point 4 (Minor): Introduction Ends Abruptly}

\begin{reviewercomment}
The introduction stops halfway through. Its last paragraph simply reads ``In this paper,''.
\end{reviewercomment}

\noindent\textbf{Location in Revised Manuscript:} Section~1 (Introduction), closing paragraph.\\[4pt]
\textbf{Response:} I have completed the introduction with a closing paragraph that previews the paper's theoretical review, data and methods, results, and discussion sections.

\subsection*{Point 5 (Minor): Definition of ``Excessive'' Attachment}

\begin{reviewercomment}
The authors talk about `excessive' attachment to the nation. What are we talking about here? What would be `healthy' attachment'? What is the reference point?
\end{reviewercomment}

\noindent\textbf{Location in Revised Manuscript:} Section~2.1 (The Needs-Based and Irrational Patriotism Models), first paragraph.\\[4pt]
\textbf{Response:} The reviewer is correct that ``excessive'' attachment implicitly invokes a reference point---some presumed healthy or normal quantity of attachment---that the original submission never specified, and that the term does not actually correspond to how the authoritarian-personality and blind-patriotism literatures define the construct I had in mind. I have removed ``excessive'' from the manuscript. In its place, I now draw explicitly on \citet{schatz1999varieties}'s distinction between \textit{blind} and \textit{constructive} patriotism, which defines the problematic form of national attachment not by its quantity but by a specific qualitative configuration: rigid, unconditional allegiance, intolerance of criticism of the nation, and uncritical endorsement of its actions, as against a constructive patriotism that remains attached to country while staying open to internal criticism aimed at reform. This reframing removes the implicit (and previously unstated) comparison to a ``healthy'' amount of attachment and instead ties the construct to an established, independently validated measure already used elsewhere in the paper (it underlies my Nationalism scale, see Section~3.2), so that ``blind'' patriotism is defined by the same criteria throughout the manuscript rather than by an undefined adjective in the theory section alone.

\clearpage

\section*{Response to Reviewer 2}

\subsection*{Point 1: Relationship Between Needs-Based Theory and Social Identity Theory}

\begin{reviewercomment}
The needs-based approach is tied primarily to the authoritarian personality tradition in the paper, and then a footnote on social identity theory and self-categorization theory claims it also is a needs-based theory because of its core self-enhancement motivation. The author is not explicit about the connection between the authoritarian personality tradition and social identity theory. The footnote also muddies the waters a bit by references to the sense of efficacy. This raises a question about how distinct are the two theoretical approaches because the importance of a ``sense of control'' is a premise of the choice process theory. A solution, in my view, is to explicitly emphasize the positive self-esteem/self-enhancement motivation that is foundational in social identity theory. This seems to be justified by the theory, and it also most clearly fits the ``needs based'' aspect of authoritarian personality. Finally, the general content of the footnote should be integrated into the text.
\end{reviewercomment}

\noindent\textbf{Location in Revised Manuscript:} Section~2.1 (The Needs-Based and Irrational Patriotism Models), second paragraph.\\[4pt]
\textbf{Response:} I agree with all three components of this comment and have revised accordingly. First, the discussion of Social Identity Theory and Self-Categorization Theory has been moved out of a footnote and into the main text, immediately following the discussion of the authoritarian-personality tradition, so that the two needs-based accounts are presented as a continuous, connected discussion rather than as a main argument plus an isolated aside.

Second, I have made the connection between the two traditions explicit rather than asserted. The revised text states directly that social identity accounts locate the source of national attachment in a self-esteem/self-enhancement motive---individuals derive positive self-regard from membership in groups perceived as prestigious, and are consequently motivated to identify with, and express pride in, high-status collectivities such as the nation---and that, although this is a different specific psychological need than the order/certainty/ego-defense motive emphasized by the authoritarian-personality tradition, it is a needs-based explanation in the same structural sense: both traditions explain national attachment as servicing a standing motivational deficit that exists prior to, and independently of, any particular structural position or lived experience of autonomy. I identify this shared structural feature explicitly as what places the two traditions together on the needs-based side of the ledger, notwithstanding their substantive differences.

Third, per the reviewer's specific suggestion, I replaced the earlier reference to a sense of ``efficacy'' with the positive self-esteem/self-enhancement motivation that is foundational to Social Identity Theory. This was the source of the conflation the reviewer flagged: ``efficacy,'' phrased imprecisely, sounded too similar to the ``sense of control'' that is the defining premise of the choice-process theory, obscuring the distinction between the two approaches. The revised paragraph now closes by stating this distinction directly: unlike the needs-based accounts, the choice-process theory does not treat affective attachment to the nation as a means of servicing a standing need for positive self-regard; attachment arises specifically out of the sense of control generated by autonomous action in concrete social contexts, and is only secondarily reattributed to the national collectivity when structural conditions make that reattribution salient. I believe this revision makes clear that the two approaches share a needs-based logic while resting on different underlying motives, and that this is precisely what differentiates them, in turn, from the control-based mechanism at the center of the choice-process theory.

\subsection*{Point 2: Introduce the Proximal/Distal Distinction Earlier}

\begin{reviewercomment}
In the treatment of choice process theory, the paper needs to introduce the distinction between ties to proximal and distal groups earlier in the paper. The choice process theory predicts stronger ties to proximal groups and treats ties to larger groups as problematic. How and when these ties spread upward to larger groups is not completely resolved in the 1992 paper. This paper suggests that a serious threat to the larger distal group is one answer to this question and finds evidence to support this claim. However, this subgroup focus of the theory needs to be explicated more clearly early in the paper, rather than emerge later as is currently the case. In this context, the results for subgroup embeddedness (membership in voluntary associations) are an important contribution of the paper. Such results offer clear support the idea that subgroup (proximal) ties tend to foster stronger distal (larger group) ties when the larger group faces a serious threat or crisis.
\end{reviewercomment}

\noindent\textbf{Location in Revised Manuscript:} Section~1 (Introduction), third paragraph.\\[4pt]
\textbf{Response:} I agree that the proximal/distal distinction was introduced too late in the original submission. The Introduction now defines the terms directly, ahead of the fuller theoretical treatment in Section~2.2: a \textit{proximal rule} ordinarily governs which of a person's nested group memberships elicits the stronger affective attachment, with local, proximal groups (within which autonomy and freedom of action are directly and repeatedly experienced) typically commanding greater loyalty than the more distant, \textit{distal} groups within which they are nested. I then state the paper's central theoretical question in these terms in the same paragraph: how, and under what conditions, this default proximal rule can reverse into a distal rule, such that affective attachment to the encompassing, distal collectivity (the nation) comes to rival attachment to the local, proximal groups nested within it. This lets the reader carry the proximal/distal vocabulary, and the paper's organizing question, from the opening pages onward, rather than encountering it only once the full theory section is reached. I agree with the reviewer that the voluntary-association membership results are an important piece of evidence for how subgroup (proximal) ties can foster stronger distal ties under threat, and I now flag this connection explicitly where those results are interpreted in the Discussion.

\subsection*{Point 3: Improve Flow and Prose}

\begin{reviewercomment}
The paper could use a careful re-write to improve the flow and clarify some of the key points. The writing is OK but not great. It is rough in places and there are long sentences that can be a bit confusing. My impression -- right or wrong -- is that the author may have rushed the final stages of preparing the manuscript. Another way to say this is: Try give the paper a bit more polish.
\end{reviewercomment}

\noindent\textbf{Location in Revised Manuscript:} Throughout.\\[4pt]
\textbf{Response:} I have revised the prose throughout the manuscript for flow and clarity. In the course of addressing the editors' and reviewers' theoretical and empirical comments above, I rewrote the affected passages---the theory sections (2.1--2.3), the introduction of the proximal/distal distinction, the Data and Variables section, the Analytic Strategy subsection, and the Discussion---breaking up several of the longer, more convoluted sentences from the original submission and tightening the exposition of the choice-process mechanism. I have additionally gone through the remainder of the manuscript with an eye toward the same issue, shortening overly long sentences and clarifying transitions between paragraphs.

\subsection*{Point 4: Supplemental Models on Political Ideology}

\begin{reviewercomment}
Delete the supplemental models on political ideology. It doesn't add much and it isn't really that necessary. It is too much of a side trip. If really important to the author, a brief footnote and offer to send the results to the reader upon request should be sufficient.
\end{reviewercomment}

\noindent\textbf{Location in Revised Manuscript:} Section~3.4 (Data and Variables), footnote; Discussion, Section~5.3 (Limitations).\\[4pt]
\textbf{Response:} I have removed Appendix A entirely, including the associated political-ideology and 1996-subsample robustness tables and all in-text references to it. In its place, Section~3.4 now carries a footnote explaining why political ideology could not be included in the primary models (in the 2004 GSS wave, the political-ideology and national-attachment modules were administered on mutually exclusive ballots) and briefly summarizing, in a single sentence, that a supplementary check on the 1996 subsample with valid ideology data showed the core education effects to be robust to a conservatism control; the footnote states that full results are available from the author upon request rather than being reported in the paper. I similarly condensed the corresponding discussion in the Limitations subsection (Section~5.3) to a brief mention rather than a full presentation of the supplementary models.

\clearpage

\section*{Summary of Comments Not Addressed}

Every substantive point raised by the editors and the two reviewers has been adopted in some form in the revised manuscript; there are no comments that I have declined to address. I do, however, flag one point on which my response is a partial, rather than a complete, resolution, and one point on which I made an editorial judgment call that reviewers may wish to revisit.

\textit{Reviewer~1, Point~3(a) (confounding events between 1996 and 2004):} As detailed above, I have reframed the paper's object of interest as the broader post-9/11 period rather than 9/11 as an isolated stimulus, and I explicitly acknowledge in the Limitations subsection (Section~5.3) that my two-wave design cannot cleanly decompose the unique contribution of the September 11 attacks from the subsequent wartime mobilization that followed. This is a conceptual reframing and an honest disclosure of a design limitation, not a resolution of the underlying identification problem; doing so would require repeated cross-sections spanning 2001--2004, which do not exist in the GSS national-attachment module. I think this is the correct way to handle the concern given the available data, but I note it here rather than implying the issue has been fully resolved.

\textit{Reviewer~2, Point~4 (supplemental political-ideology models):} The reviewer's suggestion---delete the supplemental models, or retain only a brief footnote offering results on request---explicitly offered me a choice. I opted for the footnote-only route (removing Appendix A entirely) rather than retaining an abbreviated appendix, on the judgment that this better serves the paper's length and focus; the full supplementary results remain available from the author on request, as the reviewer suggested; I note the decision here in case the reviewer's preference was for some intermediate treatment.

\vspace{12pt}
\noindent I once again thank the Editors and the two Reviewers for their time and constructive guidance, which have helped me produce a substantially stronger manuscript.

\bibliographystyle{apalike}
\bibliography{references}

\end{document}
